\section{Архитектура решения}
\label{sec:arch}

\subsection{Два этапа}
Решение состоит из предобработки (один раз) и итераций (много раз).
\begin{itemize}[nosep]
  \item \textbf{Предобработка} переводит входной CSV в компактный бинарный вид, удобный для
        потокового чтения. Это разовая плата, которая окупается: бинарный файл меньше CSV и упорядочен
        так, что итерация сводится к последовательному чтению.
  \item \textbf{Итерация} --- это шаг степенного метода LeaderRank~\eqref{eq:pull}. Бинарный файл
        читается потоком, считается ядро. В памяти живут только векторы $\bigO(n)$.
\end{itemize}
После сходимости плотные идентификаторы переводятся обратно в исходные \texttt{int32}, и результат
пишется в выходной CSV.

\subsection{Формат хранения}
\label{sec:format}
Ядру нужно для каждого приёмника $i$ перечислить его входящих соседей $j$. То есть нужна
транспонированная матрица $A\T$ --- граф, сгруппированный по приёмнику. Храним его как разреженный
столбцовый формат (CSC), разделённый на две части по принципу <<что влезает в ОЗУ --- в ОЗУ, что нет
--- на диск>>.
\begin{itemize}[itemsep=2pt]
  \item \texttt{sourcesPtr} --- массив \texttt{int[n+1]} указателей. Входящие рёбра приёмника $i$
        лежат в файле от позиции \texttt{sourcesPtr}$[i]$ до \texttt{sourcesPtr}$[i{+}1]$. Размер
        $\bigO(n)$, в ОЗУ.
  \item \texttt{sources} --- поток источников $j$, сгруппированных по приёмнику. Размер $\bigO(m)$,
        это и есть <<гигабайт>>, лежит на диске и читается потоком. Каждое ребро --- это $4$-байтовое
        целое (little-endian), только источник. Приёмник восстанавливается по положению в файле через
        \texttt{sourcesPtr}, поэтому повторно он не хранится.
  \item \texttt{outdeg} --- массив \texttt{int[n]} исходящих степеней (нужен для знаменателя
        $\degout(j)+1$). В ОЗУ.
  \item \texttt{originalIds} --- соответствие <<плотный id $\to$ исходный \texttt{int32}>> для вывода.
        В ОЗУ, нужен только при записи результата.
\end{itemize}

\subsection{Конвейер предобработки}
\label{sec:preprocess}
Предобработка --- это по сути внешняя сортировка: рёбра перегруппировываются по приёмнику, не
загружая их все в память. Именно эту <<работу с памятью>> условие требует написать самим. Этапы:
\begin{enumerate}[itemsep=3pt]
  \item \textbf{Один проход по CSV: переименование id и подсчёт степеней.} Исходные \texttt{int32}
        разрежены, поэтому они переводятся в плотный $0..n{-}1$ хеш-таблицей
        <<\texttt{int}$\to$\texttt{int}>> на примитивах. Таблица размера $\bigO(n)$ (вершины влезают в
        память, в отличие от рёбер). Заодно считаются входящие степени (для \texttt{sourcesPtr}) и
        исходящие (для знаменателя). Порядок
        присвоения плотных id фиксирован, по первому появлению. Это часть детерминизма. Чтобы не
        разбирать CSV второй раз, на этом же проходе рёбра в плотных id пишутся во временный бинарный
        файл.
  \item \textbf{Планирование корзин.} Диапазон приёмников делят на корзины так, чтобы в каждую
        попадало не больше $M$ рёбер. Размер $M$ подбирается под доступную память. Если у одной
        вершины входящая степень больше $M$, она не делится и попадает в отдельную <<большую>> корзину.
  \item \textbf{Раздача рёбер по корзинам.} Временный бинарный файл читается, каждое ребро
        дописывается в файл своей корзины. Записи последовательные, буферы фиксированного размера.
        Если корзин больше, чем можно открыть файлов под лимитом памяти, раздача идёт волнами:
        источник перечитывается несколько раз, по группе корзин за проход.
  \item \textbf{Сортировка корзин и сборка файла \texttt{sources}.} Обычная корзина целиком
        загружается в память, сортируется по ключу <<(приёмник, источник)>> и дописывается в итоговый
        файл. Большая корзина (гипер-узел) сортируется внешней сортировкой слиянием: режется на куски
        по $M$, каждый кусок сортируется и сливается потоком. Так память на одну корзину остаётся
        $\bigO(M)$ при любой степени вершины.
\end{enumerate}

Итог предобработки --- файл \texttt{sources} на диске плюс массивы $\bigO(n)$ в памяти. Дальше
итерация работает только с ними.

\paragraph{Почему детерминированно.} Итоговая раскладка (по приёмнику, источники по возрастанию) ---
это чистая функция от графа. Она не зависит ни от $M$, ни от границ корзин, ни от числа потоков. Это
проверяется тестом: $M=\infty$ и $M=$\,маленькое дают байт в байт одинаковый файл.
